% !TEX root = ../main.tex
\chapter{Conclusiones y Trabajo Futuro}
\label{chap:conclusiones}

Este capítulo final sintetiza las contribuciones principales de este Trabajo de Fin de Grado, discute sus implicaciones para el sector \textit{retail} y propone líneas de investigación para continuar evolucionando el sistema.

\section{Conclusiones Generales}
El desarrollo de este sistema ha ido más allá de un ejercicio de \textit{Machine Learning}: ha producido un \textbf{Sistema de Apoyo a la Decisión (DSS)} que conecta el modelo predictivo con la política de reposición. El trabajo valida empíricamente que la calidad de un modelo de previsión no puede medirse solo por el error puntual (MAE), sino por su capacidad para inducir decisiones logísticas económicamente mejores bajo incertidumbre.

Antes de desarrollar las conclusiones temáticas, se revisa el grado de consecución de cada uno de los objetivos específicos planteados en la Sección~\ref{sec:objetivos}, cuya verificación cuantitativa detallada se recoge en la Tabla~\ref{tab:kpi_verification}:

\begin{itemize}
    \item \textbf{\ref{obj:pipeline} (Pipeline de datos e imputación de censura): alcanzado.} El pipeline reproducible quedó operativo (Capítulos~\ref{chap:analisis_datos} y~\ref{chap:implementacion}), la ausencia de fuga de información se verificó mediante el contrato causal de las variables, y la imputación de demanda latente demostró su valor económico: $-92{,}15\%$ de coste bajo evaluación justa sobre el panel de 500 series.
    \item \textbf{\ref{obj:modelos} (Modelos globales probabilísticos): alcanzado.} El modelo global probabilístico optimizado sobre la capa conformal reduce el coste logístico un $-9{,}03\%$ frente al \textit{Seasonal Naïve} ($147.646$ frente a $162.299$ u.m.), elevando drásticamente el nivel de servicio del $52{,}45\%$ al $83{,}55\%$ ($+31{,}10$ puntos de demanda atendida).
    \item \textbf{\ref{obj:conformal} (Calibración conformal $\geq 80\%$): alcanzado.} La cobertura empírica alcanza el 84{,}6\% en CatBoost y el 85{,}5\% en LightGBM en el panel de 500 series, cumpliendo la garantía nominal del $80\%$ tras la optimización desacoplada de hiperparámetros sobre la capa conformal.
    \item \textbf{\ref{obj:mlops} (Operativización MLOps): alcanzado.} El sistema quedó desplegado como aplicación Django que sirve tanto el cuadro de mandos como la API operacional, con trazabilidad por artefactos versionados (\textit{hash} de configuración y \textit{commit} de Git en cada ejecución), explicabilidad SHAP, contenedores Docker e integración continua (Capítulo~\ref{chap:implementacion}).
\end{itemize}

De los cuatro objetivos específicos, todos se cumplen íntegramente. Las conclusiones principales se resumen en los siguientes puntos:

\begin{itemize}
    \item \textbf{Censura y Demanda Latente}: La imputación de la demanda censurada por \textit{stockouts} es el factor técnico más determinante para la rentabilidad de la cadena de suministro. Bajo una evaluación justa con un target común de demanda real, los modelos entrenados con la estrategia de Demanda Latente Supervisada reducen el coste total de inventario en un 92{,}15\% frente al baseline observado sin corregir, elevando la disponibilidad volumétrica (\textit{fill rate}) del 60{,}30\% al 97{,}50\%. Este resultado valida de forma cuantitativa y sin sesgo de evaluación el beneficio de recuperar la señal latente, con el alcance acotado al panel de 500 series del que se muestrea el protocolo (Sección~\ref{sec:resultados_gemelo}).

    \item \textbf{La trampa del MAE y el valor del modelo probabilístico}: El \textit{Seasonal Naïve} obtiene el mejor MAE del backtesting predictivo (\textbf{6,61} bajo estrategia Latent), y sin embargo a escala de 500 series ese menor error puntual es una trampa matemática: sufre roturas en casi la mitad de los días (nivel de servicio del $52{,}45\%$) y acumula $123.858$ u.m. en penalizaciones por rotura. En cambio, los modelos probabilísticos de \textit{Gradient Boosting} elevan el nivel de servicio al $83{,}55\%$ ($+31{,}10$ puntos) y reducen el coste logístico total a $147.646$ u.m. Este resultado reproduce el fenómeno teórico advertido en el Capítulo~\ref{chap:estado_del_arte}: optimizar el MAE no equivale a optimizar la política de inventario.

    \item \textbf{Calibración Conformal a Escala}: La implementación de \textit{Mondrian Conformal Prediction} proporciona una calibración condicional por categoría sólida y estable. A escala de 500 series, CatBoost alcanza una cobertura del \textbf{84{,}6\%} y LightGBM un \textbf{85{,}5\%} (ambos por encima del $80\%$ nominal), logrando LightGBM el mejor Winkler Score del sistema ($36{,}11$). La optimización desacoplada de hiperparámetros a través de la capa conformal resuelve definitivamente las desviaciones de escala observadas en conjuntos de calibración reducidos.

    \item \textbf{Los dos ejes de la evaluación económica}: el coste y el nivel de servicio no ordenan igual a los modelos, y leer solo uno de los dos induce a error. A 500 series el coste separa a los tres modelos por menos de un 2\%, mientras que el nivel de servicio los separa por 15 puntos y en el orden inverso al del MAE. La validación a escala reveló además que el beneficio de la corrección de censura es sensible a la composición del catálogo: se concentra en series de stockout severo y se diluye cuando el catálogo es más heterogéneo.


    \item \textbf{Madurez MLOps y Operativización Industrial}: El despliegue del sistema como plataforma integral (\textbf{Django} y API REST operacional), respaldada por el registro de experimentos en MLflow, promoción estadística del modelo campeón mediante guardarraíles, monitorización continua de deriva (PSI y Page-Hinkley) y una suite de tests que verifica formalmente la ausencia de fuga de información temporal, demuestra la viabilidad de transformar algoritmos probabilísticos complejos en un producto de software robusto, reproducible y listo para integrarse con los sistemas transaccionales (ERP) de una organización. La arquitectura mantiene una separación estricta entre la lógica matemática de dominio y la capa de servicio, permitiendo auditar y verificar cualquier decisión logística de forma totalmente desacoplada de la interfaz.
\end{itemize}

\section{Legado del Trabajo}
\label{sec:legado}

Más allá de las conclusiones metodológicas, este TFG deja como legado un sistema completo, verificable y reproducible:

\begin{itemize}
    \item \textbf{Código y reproducibilidad}: la totalidad del código (pipeline de datos, motor de modelos, calibración conformal, decisión Newsvendor, API REST y cuadro de mandos) se encuentra versionada en un repositorio Git público\footnote{\url{https://github.com/ArtemMindlin/retail-demand-forecasting-system}}, con integración continua (verificación de estilo, tipos estáticos y una suite de tests con contratos formales de no fuga de información temporal en cada \textit{commit}). La arquitectura completa se despliega con un único comando mediante \texttt{docker-compose}, permitiendo reproducir los experimentos y obtener resultados idénticos sin configuración manual. Cada ejecución persiste además su configuración YAML, métricas, predicciones y el \textit{hash} de Git asociado bajo el directorio \texttt{reports/} y en el servidor de seguimiento MLflow.
    \item \textbf{Datos}: el trabajo se apoya íntegramente en el dataset público \textit{FreshRetailNet-50K} \cite{freshretailnet2025}, disponible en abierto bajo licencia CC BY 4.0, por lo que la reproducción del análisis no depende de datos propietarios ni de acuerdos de confidencialidad.
    \item \textbf{Elemento persistente y demostrador}: el sistema está desplegado como instancia funcional en un servidor propio, sirviendo tanto el cuadro de mandos interactivo como la API REST operacional (documentada bajo OpenAPI/Swagger) sobre los modelos entrenados. A corto plazo sirve de soporte interactivo para la defensa; a medio plazo, la contenerización facilita su despliegue en infraestructuras corporativas; y a largo plazo, las líneas de la Sección~\ref{sec:trabajo_futuro} marcan su evolución hacia un entorno productivo industrial.
    \item \textbf{Transferencia y beneficiarios}: los responsables de aprovisionamiento de \textit{retail} disponen de una metodología cuantitativa integral (imputación de censura, previsión probabilística y \textit{Newsvendor} con restricciones de capacidad) directamente transferible a sus datos; y la comunidad académica cuenta con una implementación de referencia abierta que conecta \textit{Conformal Prediction} con la toma de decisiones económicas bajo incertidumbre.
\end{itemize}

\section{Sostenibilidad e Impacto Ambiental}
\label{sec:sostenibilidad}

El trabajo presenta una contribución directa y cuantificada a la sostenibilidad a través de la optimización de la cadena de suministro de alimentos perecederos:

\begin{itemize}
    \item \textbf{ODS 12 (Producción y Consumo Responsables - Vinculación Alta):} Constituye el núcleo del impacto ambiental del proyecto. El desperdicio alimentario en la distribución europea genera costes millonarios y pérdidas masivas de biomasa. La reducción empírica del \textbf{$92{,}15\%$} en el coste de inventario obtenida mediante la corrección de demanda latente supervisada y la calibración conformal se traduce directamente en una drástica disminución del sobrestock, evitando que toneladas de producto fresco caduquen en tienda o almacén.
    \item \textbf{ODS 13 (Acción por el Clima - Vinculación Media):} Actúa como externalidad positiva indirecta: ajustar los pedidos a la demanda real reduce la sobreproducción agrícola en origen, el consumo energético en transporte refrigerado y la huella de carbono asociada a la logística inversa de destrucción de excedentes.
    \item \textbf{ODS 2 y ODS 9 (Hambre Cero e Innovación):} La elevación de la disponibilidad de producto ($97{,}50\%$ de \textit{fill rate}) mejora el abastecimiento minorista, mientras que el desarrollo de un pipeline MLOps abierto promueve infraestructuras tecnológicas reproducibles.
\end{itemize}

El análisis detallado y la matriz oficial de vinculación con los 17 ODS de la Agenda 2030 se recogen en el Anexo~\ref{anx:ods}, conforme a la normativa de la ETSINF-UPV.

\section{Análisis de Impacto Ético, Social y Privacidad}
\label{sec:impacto_etico}

Más allá de los resultados logísticos, el sistema se ha estructurado atendiendo a criterios de ética algorítmica, responsabilidad social corporativa y el marco regulatorio europeo:

\begin{itemize}
    \item \textbf{Ética de la Automatización y Supervisión Humana (\textit{Human-in-the-Loop})}: En consonancia con las directrices de IA Confiable de la Comisión Europea \cite{eucommission2019aiethics} y el Reglamento Europeo de Inteligencia Artificial (\textit{EU AI Act}) \cite{euaiact2024}, el sistema se concibe estrictamente como una herramienta de apoyo a la decisión (DSS) y no como un mecanismo de sustitución ciega del criterio humano. El uso de \textbf{SHAP} para auditar la contribución de cada variable, junto al panel de gestión por excepciones que señala anomalías (deriva temporal, \textit{cold start} o pedidos atípicos), garantiza una colaboración humano-máquina transparente, donde el operador mantiene siempre el control final.

    \item \textbf{Consumo Energético y Eficiencia Computacional}: El entrenamiento de modelos globales sobre millones de registros requiere recursos computacionales no despreciables. No obstante, la arquitectura MLOps desacoplada (búsqueda aperiódica de hiperparámetros y reentrenamiento semanal ligero de apenas $5{,}5$ segundos) minimiza el ciclo de cómputo intensivo, resultando en una huella de carbono marginal frente al volumen de emisiones y desperdicio alimentario que evita en la cadena logística.

    \item \textbf{Privacidad y Cumplimiento RGPD}: El dataset \textit{FreshRetailNet-50K} está completamente anonimizado, satisfaciendo el Reglamento General de Protección de Datos (RGPD / GDPR). El pipeline procesa exclusivamente series agregadas de demanda física a nivel tienda-producto, sin almacenar, inferir ni tratar ningún dato de carácter personal de consumidores.

    \item \textbf{Propiedad Intelectual y Licenciamiento Abierto}: \textit{FreshRetailNet-50K} \cite{freshretailnet2025} se distribuye bajo licencia \textbf{CC BY 4.0}, cuyos términos se satisfacen mediante la debida atribución bibliográfica. El código fuente desarrollado se publica íntegramente bajo licencia permisiva \textbf{MIT} en el repositorio del proyecto, facilitando su reutilización por la comunidad.

    \item \textbf{Equidad Algorítmica y Sesgo Territorial}: El sistema no incorpora variables sociodemográficas ni atributos protegidos. No obstante, se advierte un riesgo indirecto de sesgo estructural: tiendas de reciente apertura o ubicadas en zonas con menor densidad comercial presentan historiales más cortos e intermitentes, lo que genera estratos conformales con mayor dispersión residual y, potencialmente, un nivel de servicio inferior por escasez de muestra. Reconocer esta limitación permite a los gestores monitorizar activamente estos puntos de venta periféricos y aplicar amortiguadores de inventario específicos.
\end{itemize}

\section{Limitaciones Metodológicas y del Sistema}
\label{sec:limitaciones}

El diseño e implementación del sistema presentan ciertas simplificaciones operativas, metodológicas y tecnológicas que es necesario advertir para contextualizar el alcance de los resultados:

\begin{itemize}
    \item \textbf{Decisión Newsvendor Desacoplada sin Arrastre de Estado}: La evaluación económica Newsvendor resuelve la cantidad de reposición de forma independiente en cada origen de decisión. No modela el arrastre dinámico de existencias sobrantes entre jornadas consecutivas, los tiempos de suministro estocásticos (\textit{lead time}) ni las restricciones de capacidad volumétrica de almacén o transporte.

    \item \textbf{Perfiles de Coste Sintéticos frente a Contabilidad ERP}: Dado que el dataset público no incluye márgenes comerciales ni costes reales de merma y rotura por SKU, el sistema recurre a perfiles de coste normalizados ($c_u, c_o$). Aunque preservan la estructura económica del problema, la traslación a producción real exigirá conectar el Newsvendor con los costes contables del ERP corporativo.

    \item \textbf{Granularidad Diaria en la Reconstrucción de Censura}: El módulo \texttt{Latent}\allowbreak\texttt{Demand}\allowbreak\texttt{Imputer} opera a nivel diario agregado: si una jornada registra horas de rotura, imputa la demanda del día completo. No realiza la reconstrucción quirúrgica hora a hora intra-día que permitiría el registro granular del dataset original.

    \item \textbf{Supuesto de Intercambiabilidad en Conformal Prediction}: La garantía teórica de cobertura de la predicción conforme descansa en la intercambiabilidad entre las muestras de calibración y evaluación. En horizontes temporales afectados por derivas pronunciadas de demanda o catálogos de extrema intermitencia, se produce una leve pérdida de calibración condicional que los métodos estáticos no corrigen dinámicamente.

    \item \textbf{Infraestructura de Despliegue y Seguridad Empresarial}: El sistema se encuentra actualmente desplegado como una instancia de demostración funcional accesible públicamente (mediante Docker Compose y proxy inverso), respaldada por una base de datos SQLite y almacenamiento de artefactos en volumen local. Si bien esto resulta idóneo para la validación académica y la defensa del proyecto, un entorno corporativo real exigiría confinar el acceso a la red privada corporativa (VPC/VPN) con autenticación federada (SSO/IAM), así como escalar la persistencia hacia bases de datos transaccionales distribuidas (PostgreSQL), almacenamiento de objetos en la nube (S3/GCS), orquestación en Kubernetes y \textit{Feature Stores} en tiempo real.
\end{itemize}

\section{Análisis de Riesgos}
\label{sec:analisis_riesgos}

Las limitaciones de la sección anterior describen lo que el trabajo no cubre. Esta recoge lo distinto: qué puede salir mal si el sistema se lleva a explotación, con qué impacto y qué medida lo contiene. Se ordenan por criticidad, y en todos los casos la medida citada está implementada, no es una intención.

\begin{itemize}
    \item \textbf{Riesgo tecnológico: descalibración del imputador al cambiar de escala.} El fichero de hiperparámetros del imputador supervisado solo es válido cerca del tamaño de ajuste con el que se obtuvo, y el efecto llega a invertirse de signo (Anexo~\ref{anx:tuning_imputador}). Un despliegue sobre un catálogo de tamaño distinto degradaría la reconstrucción sin emitir error alguno, y como toda la cadena se alimenta de esa señal, el deterioro llegaría hasta la cantidad de pedido. \emph{Medida}: los metadatos de cada búsqueda registran las filas limpias de ajuste, y la resolución de parámetros rechaza un fichero cuya huella de escala no corresponda, recurriendo entonces a los valores por defecto de la configuración.

    \item \textbf{Riesgo tecnológico: pérdida de cobertura conformal por deriva.} La garantía de la predicción conforme descansa en la intercambiabilidad entre la ventana de calibración y la de test, supuesto que la deriva temporal rompe. Una infracobertura sostenida se traduce en roturas de stock por debajo del nivel de servicio comprometido. \emph{Medida}: el índice PSI sobre las variables más influyentes y el detector Page-Hinkley sobre el error vigilan las dos derivas por separado (Sección~\ref{sec:monitorizacion}), y la cobertura empírica se reporta por fold en cada ejecución en lugar de asumirse.

    \item \textbf{Riesgo de integración: modelo campeón obsoleto.} Un campeón promovido y no revisado se degrada a medida que el régimen de demanda se desplaza. \emph{Medida}: la promoción está sujeta a un criterio automático de coste sobre el registro de campeones, y la cadencia de reentrenamiento se eligió empíricamente (Sección~\ref{sec:resultados_simulacion}) en lugar de fijarse por convención.

    \item \textbf{Riesgo de aceptación: desconfianza del planificador ante una recomendación opaca.} Un sistema de apoyo a la decisión que no se explica se termina ignorando, y el beneficio medido no se materializa por mucho que el modelo acierte. \emph{Medida}: el cuadro de mandos publica la atribución SHAP del campeón y la gestión por excepciones prioriza los casos que requieren criterio humano, de modo que la herramienta actúa como copiloto y no como caja negra. Como indicador de aceptación medible tras una implantación real, la fracción de recomendaciones aceptadas sin modificación sería la métrica natural a instrumentar.

    \item \textbf{Riesgo operativo: pérdida de la trazabilidad de las corridas.} El almacén de corridas se respalda en un fichero SQLite sin réplica, de modo que su pérdida dejaría los artefactos sin la información que los identifica. \emph{Medida parcial}: cada directorio de artefactos incluye un fichero \texttt{mlflow\_run.json} con el nombre y el identificador de su corrida, lo que permite reconstruir la correspondencia; la réplica del almacén queda como requisito de un entorno productivo, según lo indicado en la Sección~\ref{sec:limitaciones}.
\end{itemize}

\section{Relación del Trabajo con los Estudios Cursados}
\label{sec:relacion_estudios}

Este trabajo articula las bases formativas adquiridas en el \textbf{Grado en Ciencia de Datos} de la Escola Tècnica Superior d'Enginyeria Informàtica (ETSINF-UPV), complementándolas con un esfuerzo sustancial de investigación y desarrollo de ingeniería autónomo:

\begin{itemize}
    \item \textbf{Análisis y preparación de datos}: La caracterización de la intermitencia, estacionalidad y censura del dataset (Capítulo~\ref{chap:analisis_datos}) aplica los fundamentos de \textit{Análisis Exploratorio de Datos}, \textit{Visualización}, \textit{Gestión de Datos}, \textit{Bases de Datos} y las asignaturas de \textit{Proyecto I y II}.

    \item \textbf{Modelización predictiva y estadística inferencial}: El ajuste de modelos globales de \textit{Gradient Boosting}, la regresión cuantílica y el marco probabilístico de decisión bajo incertidumbre que sustenta el \textit{Newsvendor} descansan sobre \textit{Modelos Descriptivos y Predictivos I y II} y \textit{Modelos Estadísticos para la Toma de Decisiones I y II}.

    \item \textbf{Metodología de validación}: El diseño del \textit{backtesting} temporal sin fuga de información y el análisis conceptual de deriva (\textit{drift}) se apoyan en \textit{Evaluación, Despliegue y Monitorización de Modelos}.

    \item \textbf{Dimensión de negocio y marco legal}: La traducción de predicciones a coste logístico, el presupuesto económico (Capítulo~\ref{chap:presupuesto}) y el análisis de impacto ético y RGPD (Sección~\ref{sec:impacto_etico}) aplican competencias de \textit{Fundamentos de Organización de Empresas}, \textit{Economía Digital} y \textit{Marco Profesional, Legal y Deontológico}.
\end{itemize}

\textbf{Aprendizaje autónomo e ingeniería de software}: Dado que el plan de estudios se orienta prioritariamente al modelado analítico, una parte determinante del proyecto exigió un aprendizaje autónomo avanzado más allá de las aulas:
\begin{enumerate}
    \item \textit{Conformal Prediction}: El estudio e implementación de la calibración conformal (y en particular su variante estratificada \textit{Mondrian}) para obtener intervalos con cobertura garantizada.
    \item \textit{Tratamiento de la demanda censurada}: La asimilación de la literatura de imputación de demanda latente en retail para revertir el sesgo de rotura de stock.
    \item \textit{Arquitectura web, APIs y servidores}: El desarrollo de la plataforma completa en \textbf{Django} con componentes dinámicos en \textbf{HTMX}, la especificación y documentación de la \textbf{API REST con OpenAPI/Swagger}, el empaquetado en contenedores \textbf{Docker Compose}, la configuración del servidor de producción (Nginx/Gunicorn), la suite de tests automatizados con contratos temporales en \textbf{pytest} y la integración del servidor de trazabilidad \textbf{MLflow}.
\end{enumerate}

En cuanto a las \textbf{competencias transversales institucionales de la UPV}, el trabajo pone en práctica:
\begin{itemize}
    \item \textbf{CT-01 (Comprensión e integración):} Coordinación de estadística matemática, algoritmos de aprendizaje, ingeniería web y logística en un único DSS.
    \item \textbf{CT-03 (Análisis y resolución de problemas):} Descomposición modular y justificada del problema (imputación $\rightarrow$ predicción $\rightarrow$ calibración $\rightarrow$ decisión Newsvendor $\rightarrow$ simulación).
    \item \textbf{CT-07 (Responsabilidad ética, medioambiental y profesional):} Análisis de sostenibilidad (ODS 12 y 13), estricto cumplimiento del RGPD, directrices del \textit{EU AI Act} y código abierto bajo licencia MIT.
    \item \textbf{CT-09 (Pensamiento crítico):} Verificación empírica honesta, demostración de la insuficiencia del MAE frente al coste económico, diagnóstico del colapso conformal a escala reducida (Sección~\ref{sec:resultados_escala}) y declaración deliberada de no-conclusión en políticas de reentrenamiento por muestra insuficiente (Sección~\ref{sec:resultados_simulacion}).
    \item \textbf{CT-10 (Aprendizaje permanente):} Capacidad para investigar e integrar tecnologías de frontera no impartidas en el grado.
    \item \textbf{CT-11 (Planificación y gestión del tiempo):} Ejecución metódica del proyecto dentro de las 300 horas estimadas en el presupuesto.
\end{itemize}

\section{Líneas de Trabajo Futuro}
\label{sec:trabajo_futuro}
La arquitectura modular desarrollada deja el camino preparado para futuras expansiones. Las líneas propuestas fluyen directamente de las limitaciones identificadas en el capítulo de resultados y en la Sección~\ref{sec:limitaciones}:

\begin{itemize}
    \item \textbf{Conformal Prediction Adaptativa (EnbPI)}: La validación a escala demostró que,
    eliminada la inestabilidad por estratos pequeños, la brecha de cobertura restante proviene de la
    no-intercambiabilidad temporal entre calibración y test, cuyo efecto sobre la garantía está
    teóricamente acotado \cite{barber2023beyond}. La línea prioritaria es sustituir el
    Split CP estático por métodos adaptativos como EnbPI \cite{xu2021conformal} o ACI
    \cite{gibbs2021adaptive}, que reajustan respectivamente el ancho del intervalo y el nivel de
    significancia con los errores recientes observados, siguiendo el protocolo comparativo de
    Zaffran et al.\ \cite{zaffran2022adaptive}, y verificar la convergencia hacia el
    80\% nominal sobre el dataset completo (50.000 series) con historia más larga.

    \item \textbf{Costes Reales e Integración ERP}: Los perfiles de coste sintéticos ($c_u$, $c_o$)
    utilizados en este trabajo son una aproximación necesaria dado que el dataset no incluye
    márgenes reales. La integración con el ERP permitiría sustituirlos por costes empíricos de
    merma, rotura y liquidación, haciendo el Newsvendor completamente operativo sobre costes
    medidos en lugar de sobre los perfiles sintéticos de la Sección~\ref{sec:cost_profiles}.

    \item \textbf{Intervalo de Promoción Directo Frente al Campeón}: El tercer guardarraíl
    de la Sección~\ref{sec:promocion} certifica que el \textit{challenger} bate al \textit{naive}
    estacional de forma no atribuible al azar, no que bate al campeón vigente, que es la
    comparación que en última instancia decide un reemplazo de modelo en producción. La
    extensión natural es aplicar el mismo \textit{bootstrap} emparejado por conglomerados de
    serie directamente sobre el hueco \textit{challenger} contra campeón, cerrando la distancia
    entre lo que el guardarraíl mide hoy y lo que la decisión de negocio necesita.

    \item \textbf{Gemelo Digital de Inventario con Estado}: la evaluación económica actual
    resuelve una decisión independiente por origen de predicción, sin arrastre de existencias
    entre periodos. La extensión natural, y la línea prioritaria en el plano de la decisión, es un
    simulador de reposición con estado que encadene existencias, pedidos en tránsito y demanda
    insatisfecha, permitiendo medir efectos que una decisión aislada no puede exponer, como la
    amplificación de la variabilidad de los pedidos a lo largo de la cadena (\textit{efecto
    látigo}). Sobre esa base se apoyarían dos extensiones inmediatas: modelar el \textit{lead
    time} como variable aleatoria a partir del histórico de entregas, e incorporar un techo de
    capacidad de almacén que obligue a repartir el espacio entre productos cuando la suma de las
    cantidades individualmente óptimas no cabe.

    \item \textbf{Modelos de Deep Learning para Series Temporales}: Explorar arquitecturas como el
    Temporal Fusion Transformer para sustituir el motor de \textit{Gradient Boosting}, capturando
    dependencias temporales a largo plazo sin ingeniería manual de \textit{lags}. El trabajo
    de Salinas et al.\ \cite{salinas2020deepar} y los resultados de la competición M5
    \cite{makridakis2022m5} sugieren que estos modelos son competitivos cuando el volumen de
    datos es suficiente.

    \item \textbf{Recuperación de Demanda Latente a Granularidad Horaria}: El módulo
    \texttt{Latent}\allowbreak\texttt{Demand}\allowbreak\texttt{Imputer} opera a nivel diario: si un día presenta \texttt{stockout\_hours} $>
    0$, la imputación se aplica sobre la demanda agregada de ese día completo, independientemente de
    si el stockout duró 2 horas o 22. El dataset FreshRetailNet-50K contiene registros de ventas y
    estado de stock hora a hora, granularidad que el paper de referencia del dataset
    \cite{freshretailnet2025} explota explícitamente: su pipeline reconstruye la demanda latente
    únicamente en las horas sin stock y suma la demanda observada en las horas con stock, logrando
    una imputación quirúrgica que elimina el sesgo sistemático de -7{,}37\% en la predicción. La
    extensión natural de este trabajo es operar el imputer a granularidad horaria, agregar después a
    nivel diario, y comparar el efecto sobre el Winkler Score frente a la imputación diaria actual.

    \item \textbf{Imputación de Censura Multi-Producto}: El módulo actual de \texttt{Latent}\allowbreak\texttt{Demand}\allowbreak\texttt{Imputer}
    trata cada serie de forma independiente. Una extensión natural es modelar la canibalización y
    sustitución entre productos de la misma categoría durante un stockout, incorporando efectos
    cruzados en la estimación de la demanda latente.

    \item \textbf{Modelos Hurdle para Demanda Intermitente}: De las 200.760 observaciones con venta
    cero identificadas en el Capítulo~\ref{chap:analisis_datos} (un 4{,}5\% del panel), el subconjunto
    con demanda \emph{genuinamente} nula, es decir, excluyendo los ceros inducidos por stockout,
    representa el 3{,}3\% del panel. Esa fracción no justifica
    un modelo de dos etapas para el catálogo general, pero sí para los SKUs de alta intermitencia
    donde la probabilidad de venta cero es estructuralmente distinta de la magnitud cuando sí se
    vende. La arquitectura natural es un clasificador binario (si hay o no venta) seguido de un regresor
    sobre el subconjunto positivo, análogo al modelo \textit{zero-inflated} clásico
    \cite{lambert1992zero}: la primera etapa filtra los días sin demanda y la segunda predice la
    cantidad condicionada a que haya venta. En el contexto de forecasting de series temporales,
    este enfoque se ha aplicado con éxito a demanda esporádica en retail \cite{syntetos2005accuracy}.

    \item \textbf{Aprendizaje por Refuerzo para Política de Reposición}: Transicionar de un
    pipeline secuencial (pronóstico $\rightarrow$ Newsvendor) a un agente que aprenda
    directamente la política óptima de reposición interactuando con un simulador de inventario con
    estado, eliminando la necesidad de especificar los perfiles de coste de forma explícita. Esta
    línea presupone el gemelo digital descrito más arriba.
\end{itemize}

Como cierre, este TFG demuestra empíricamente que el valor de la inteligencia artificial aplicada a operaciones no está solo en los algoritmos, sino en integrarlos dentro de un diseño de software que cierra el bucle desde el dato bruto hasta la decisión logística operativa, con sus garantías, sus limitaciones cuantificadas y su trazabilidad completa.
